% Manuscript for Radiology: Artificial Intelligence (Original Research), in the cai4cai/RSNATemplate skeleton.
% Build: latexmk (xelatex + biber) in this folder; see README.md and TEMPLATE_NOTES.md.
% Every number in this document is traced to a result file in NUMBERS.md; only rigor-audit-passing numbers appear.
% Every claim about what this work IS is adjudicated in CLAIMS_LEDGER.md against the published comparator, and the
% sentence here is the one that ledger row says is defensible. The LevelCheck comparator (LEVELCHECK_MODEL.md, frozen
% by the commit that added it, ecae75f, last d0403de) belongs to paper 2 since 2026-09-24 (paper2_levelcheck/).
% A quantity a gate has not yet produced is \PENDING{gate}{what}: it prints in red and tools/check_limits.py fails;
% a figure another workflow is producing is \PENDINGFIG{path}{gate}{what}, drawn as soon as its file exists.
% WRAPPED AROUND FIVE FIGURES (Greg, 2026-09-23 named six; on 2026-09-24 the sixth, the LevelCheck comparison, moved
% to paper 2): 1 transitional phenotypes (fig:phenotypes), 2 the encoder's embedding (fig:embedding), 3 the network
% (fig:network), 4 one exposure on an unseen C-arm (fig:exposure), 5 persons needed, the film N curve and
% acquisition-space coverage (fig:power). Figure files are SchwingNet/figures/figureN.{pdf,tif,png} with figureN.caption.txt (captions here are
% those files verbatim, LaTeX-escaped) and figureN.record.json. Tables in citation order: 1 image sources (tab:data), 2 transitional
% persons (tab:transitional), 3 naming, counting and coupling (tab:naming; file tab3_naming.tex). Supplement
% (supplement.tex): Figure S1 flow of cases (fig2_flow), S2 corpus manifold (figS1.*), S3 image-formation parameters
% (figS2.*), S4 which training source carries the detector (figS3.*), S5 detection against view angle (figS4.*);
% Table S1 device terms, Table S2 detection per device and arm. Retired from the main text: the method composite,
% the gallery, fig3_detection, fig4_formation, fig5_naming, fig6_claim, fig_manifold, fig5_contribution (their files
% stay in figures/).
% Re-wrapped 2026-09-24 against the figure builds of that morning (Figure 1 T13 example from VerSe; Figure 2 fused
% L5 and the contested never-seen L6; Figure 3 two detectors; Figure 4 L6 calls at confidence 0.74; Figure 5 the
% frozen exact test and the built split).
% Revision 2 (2026-09-24, manuscript-drive; ledger REVISION_LOG2.md): EVAL_LEAKS.md (zero-shot = the lineage-complete
% v9 encoder + head, HELD OUT; every fold-model number PENDING G5 "C-ArmLSpine CV on the v9 lineage"; the embedding
% is the selection split, said everywhere it is quoted), DRIVE_CHECK.md D1-D3, REVIEW2_RSNA.md, REVIEW2_STATS.md,
% REFS_AUDIT.md (refs_corrections.bib merged into refs.bib), LEVELCHECK_MODEL.md s14-s15. The network is named
% SchwingNet in the title, abstract, key points and text (TITLE.md, ANONYMISED-VARIANT = no). The title is fixed by
% TITLE.md; REVIEW2_RSNA Fix 4's alternative titles are Greg's decision and are not applied.
% Revision 2, round 2 (2026-09-24 10:24 onward; REVISION_LOG2.md "Round 2"): the acquisition-space review integrated
% from lit/CITABLE_CLAIMS.md (VERIFIED rows only; zheng2018pda and diffpose2024 added, rao2024 out: 35 references);
% the G5 fold-model numbers filled from results/eval_gate/closure/cv_v9/ (HELD OUT), matching the rebuilt Figure 4;
% the same-device learning curve and the arc pose pools in Methods, Results and the supplement.
% Revision 3 (2026-09-24 ~15:45 onward, collect-and-submit-1; ledger REVISION_LOG3.md): the reviews, references,
% closure section 7 and CITABLE_CLAIMS were already applied by revision 2 and were verified, not redone. New: the
% supplement section "Data integrity and quality controls" (every audit tool, one Methods sentence pointing to it);
% the Methods name the secondary analyses of LEVELCHECK_MODEL.md s12-s13 including the selective S1 (detail in
% Table 3 note c); the G5 placeholder in the Results now names the acquisition-space coverage / film-vs-fluoroscopy
% N curve (results/indomain/fit.json exists but Figure 5 has no in-domain panels and no NHANES II N curve exists);
% measured trims for room (body 2,999 of 3,000). G0, G4, G6 and the Figure 6 file remain open.
% Revision 5 (2026-09-24 evening; Greg: "move the comparison to neurosurgeons / all the LevelCheck stuff into a new
% paper"; ledger REVISION_LOG3.md section 8; scope rule PAPER1_SCOPE.md): the LevelCheck comparison (Figure 6,
% Table 3, the Methods and Results paragraphs, the Discussion paragraph, the abstract, key-point, summary-statement
% and title clauses, the Registration sentence, the power keys) moved verbatim to
% paper2_levelcheck/MOVED_FROM_PAPER1.md. Five figures, three tables, 34 references. The Introduction's prior-art
% sentences on LevelCheck stay. No transitional test person of the built split is scored here beyond the three
% C-ArmLSpine L6 cross-validation cases, which paper 2 excludes from its primary analysis in advance.
% Fill pass (2026-09-24 evening, REVISION_LOG3.md section 8, "Fill pass"): G0 dates (April--September 2026), the
% readers (no year of training recorded), age and sex through \RESULT{agesex_*} from results/age_sex_by_source.json
% (per-source detail: supplement Table S3), the Ethics line (VinDr not covered by 2026-269; notification clause
% dropped, no record it was sent), G4 census (supplement Table S4 from results/corpus_census.json), the
% automated-reference paragraph moved to paper 2 (no paper-1 result uses it). One G5 placeholder remains in the
% Results (the Figure 5 data case; grid chains 40437811, 40437814) plus Table S2's 15 G5 cells.
% Two copies from one source (Greg 2026-09-25; REVISION_LOG3.md section 10): rsnastyle.tex's \ifcameraready is FALSE by
% default (the double-anonymized REVIEW copy) and TRUE when slurm/compile.sh runs with CAMERAREADY=1 (it writes
% cameraready.flag.tex). The network's name is never typed in the text: \netname{} prints "[AuthorLastName]-Net" /
% "SchwingNet"; \withheldurl{url} prints "[URL withheld for review]" / the URL; \selfcite{key} prints "[reference
% withheld for review]" / \citep{key} (the entry then leaves the review reference list). Comments keep the name.
% tools/check_limits.py and tools/make_packet.py measure and build the copy the environment selects (TITLE.md).
% rsnastyle.tex -- preamble after the cai4cai/RSNATemplate (Radiology: Artificial Intelligence submission style).
% Reproduced from that repository as faithfully as recalled; every departure is listed in TEMPLATE_NOTES.md.
% The template assumes a full TeX Live with fontspec, biblatex + biber, endfloat, lwarp and glossaries.
% Where a package is absent (the compute grid carries TeX Live 2013 without biblatex, biber, endfloat or lwarp)
% a working minimal equivalent is selected with \IfFileExists so the same main.tex compiles in both places.
\documentclass[11pt]{article}

% ---- fonts (the template: Times New Roman through fontspec, hence xelatex or lualatex) ----------------------
\usepackage{fontspec}
\InputIfFileExists{localfont.tex}{}{\setmainfont{Times New Roman}}
% localfont.tex, when present, holds a single \setmainfont line for machines without Times New Roman
% (the grid build writes \setmainfont{Liberation Serif}, metrically compatible with Times).

% ---- HTML export hook of the template; opt-in here --------------------------------------------------------------
% The template loads lwarp unconditionally. In the PDF build it rewrites glossaries (the abbreviation list came
% out empty) and adds a blank leading page, so it is loaded only when a marker file `lwarp-enable.tex` exists.
\IfFileExists{lwarp-enable.tex}{\IfFileExists{lwarp.sty}{\usepackage{lwarp}}{}}{}

% ---- page ----------------------------------------------------------------------------------------------------------
\usepackage[a4paper,left=2.5cm,right=2.5cm,top=3cm,bottom=3cm]{geometry}
\usepackage{setspace}
\doublespacing
\usepackage{parskip}
\usepackage{graphicx}
\usepackage{booktabs}
\usepackage{array}
\usepackage{amsmath}
\usepackage{xcolor}
\usepackage[labelfont=bf,font=small]{caption}
\usepackage{authblk}
\usepackage{tocloft}

% ---- bibliography (the template: biblatex, style=nejm, natbib=true, with a simplebib enumerated environment) -----
\newif\ifrsnabiblatex
\IfFileExists{biblatex.sty}{\rsnabiblatextrue}{\rsnabiblatexfalse}
\ifrsnabiblatex
  \usepackage[style=nejm,natbib=true,backend=biber,maxbibnames=6,minbibnames=3,sorting=none]{biblatex}
  \addbibresource{refs.bib}
  % simplebib: the numbered list the template prints the references in
  \defbibenvironment{simplebib}
    {\list{\printtext[labelnumberwidth]{\printfield{labelprefix}\printfield{labelnumber}.}}
          {\setlength{\labelwidth}{\labelnumberwidth}\setlength{\leftmargin}{\labelwidth}
           \setlength{\labelsep}{\biblabelsep}\addtolength{\leftmargin}{\labelsep}
           \setlength{\itemsep}{\bibitemsep}\setlength{\parsep}{\bibparsep}}\renewcommand*{\makelabel}[1]{\hss##1}}
    {\endlist}
    {\item}
  \newcommand{\rsnabibliography}{\printbibliography[env=simplebib,title={References}]}
\else
  % fallback for a TeX install without biblatex: natbib numeric, unsorted (order of first citation), bibtex
  \usepackage[numbers,sort&compress]{natbib}
  \newcommand{\rsnabibliography}{\bibliographystyle{unsrtnat}\bibliography{refs}}
\fi

% ---- floats at the end (the template: endfloat with tablesfirst and nomarkers) -----------------------------------
\IfFileExists{endfloat.sty}{\usepackage[nomarkers,tablesfirst,nolists]{endfloat}}{}
% nolists: endfloat would otherwise print its own List of Figures / List of Tables; the figure legends are
% printed by main.tex from the lof instead, as the template does.

% ---- acronyms (the template: glossaries with the acronym option) ----------------------------------------------------
\usepackage[acronym,nonumberlist,nopostdot]{glossaries}
% the template calls \makeglossaries and \printglossary and relies on latexmkrc to run makeglossaries;
% the noidx variant (glossaries >= 4.04) needs no external program. On an older glossaries (TeX Live 2013)
% the acronyms are collected into a plain list instead, so the same main.tex compiles there too.
\makeatletter
\ifdefined\makenoidxglossaries \makenoidxglossaries \fi
% \rsnaacronym{key}{ABBR}{expansion} registers the acronym with glossaries (so \gls works) and appends it to a
% plain list that \printabbreviations prints; the printed list therefore never depends on the glossary
% machinery having run (with TeX Live 2020's glossaries 4.46 the noidx list printed empty in this document).
\def\rsna@abbr{}
\newcommand{\rsnaacronym}[3]{\ifdefined\newacronym\newacronym{#1}{#2}{#3}\fi\g@addto@macro\rsna@abbr{#2, #3; }}
\newcommand{\printabbreviations}{\rsna@abbr}
\makeatother

% ---- hyperref last (after tocloft and glossaries, which both redefine list-of-figures internals) ---------------------
\usepackage[hidelinks]{hyperref}

% ---- \PENDING: a quantity this manuscript does not have yet -----------------------------------------------------
% \PENDING{gate}{what} prints a red bold bracket that cannot be missed in a read of the PDF, and
% tools/check_limits.py exits non-zero while one remains. No number is ever invented in its place.
% \PENDINGFIG{path}{gate}{what} draws the figure when its file exists and prints the same marker when it does not,
% so a figure another workflow is producing can be wired now; check_limits.py fails only while the file is absent.
\newcommand{\PENDING}[2]{{\bfseries\color{red}[PENDING #1: #2]}}
\newcommand{\PENDINGFIG}[3]{\IfFileExists{#1}{\includegraphics[width=\linewidth]{#1}}{\PENDING{#2}{#3}}}

% ---- small helpers used in main.tex ---------------------------------------------------------------------------------
\providecommand{\ci}[2]{[#1, #2]}
\providecommand{\pct}{\,\%}
\providecommand{\degree}{$^\circ$}
\renewcommand{\degree}{$^\circ$}

% ---- two copies from one source: the double-anonymized REVIEW copy and the CAMERA-READY copy (Greg 2026-09-25) ----
% \ifcameraready is FALSE by default (review). slurm/compile.sh writes cameraready.flag.tex (\camerareadytrue) when
% it runs with CAMERAREADY=1 and removes it otherwise; nothing else sets the flag. The three macros are the only
% places the network's name, the authors' URLs and the self-citations may appear in the tex (comments excepted):
%   \netname            SchwingNet                 | [AuthorLastName]-Net
%   \withheldurl{url}   \url{url}                  | [URL withheld for review]
%   \selfcite{key}      \citep{key}                | [reference withheld for review]  (no \cite, so the entry does
%                                                    not enter the review reference list; numbering stays consecutive)
% Robust, so they survive \title, \caption (the .lof) and \textbf. tools/check_limits.py measures the copy the
% environment selects and greps the review copy for the name; tools/make_packet.py builds the same copy.
\newif\ifcameraready
\camerareadyfalse
\IfFileExists{cameraready.flag.tex}{\input{cameraready.flag.tex}}{}
\DeclareRobustCommand{\netname}{\ifcameraready SchwingNet\else [AuthorLastName]-Net\fi}
\DeclareRobustCommand{\withheldurl}[1]{\ifcameraready\url{#1}\else [URL withheld for review]\fi}
\DeclareRobustCommand{\selfcite}[1]{\ifcameraready\citep{#1}\else [reference withheld for review]\fi}

% ---- results-driven build (RESULTS_BUILD.md) ---------------------------------------------------------------
% \RESULT{key} / \RESULTFIG{key} read generated/numbers.tex, which tools/results_macros.py writes from
% numbers_manifest.json before every compile; a missing result prints in red. Loaded LAST: the package only
% \providecommand's \PENDING and \PENDINGFIG, so the definitions above stay the ones in force.
\usepackage{rsnaresults}


\title{\netname{}: a convolutional-recurrent network for level naming in spine surgery}
% Double-anonymized review: authors, affiliations, funding and acknowledgments are on the separate Full Title Page.
\author{}
\date{}

% ---- abbreviations (at most ten; CT is not listed, per the journal); alphabetical, as they print -------------------
% Revision 2: DRIVE_CHECK item 7 -- C2, L1-L4, T11 and T12 were used undefined; the level entries now cover the
% range the text uses (S1 left the text with Table 3 on 2026-09-24; F1 and JSON left the text).
\rsnaacronym{bca}{BCa}{bias-corrected and accelerated bootstrap interval}
\rsnaacronym{c2}{C2}{second cervical vertebra}
\rsnaacronym{ci}{CI}{confidence interval}
\rsnaacronym{gru}{GRU}{gated recurrent unit}
\rsnaacronym{l5}{L1--L6}{first to sixth lumbar vertebrae}
\rsnaacronym{t13}{T11--T13}{eleventh to thirteenth thoracic vertebrae}
\rsnaacronym{tsne}{t-SNE}{t-distributed stochastic neighbour embedding}
% The style file's list ends every entry with "; ", so the last one would close on a stray semicolon
% (REVIEW_RSNA.md Fix 8): print the list with the final "; " replaced by a full stop.
\makeatletter
\def\rsna@striplast#1; \@nil{#1}
\renewcommand{\printabbreviations}{\expandafter\rsna@striplast\rsna@abbr\@nil.}
\makeatother

\begin{document}
% emergency stretch: long names (CTSpinoPelvic1K, C-ArmLSpine) and DOIs otherwise overrun the line by 2-27 pt
\setlength{\emergencystretch}{2em}
\maketitle

\textbf{Article type}: Original research

% Front matter written last, from the final numbers of revision 2 (REVISION_LOG2.md). Only held-out numbers
% (EVAL_LEAKS.md CLEAN) are quoted without a qualifier; the embedding carries "the split that chose its checkpoint";
% the fold-model numbers are PENDING G5 (C-ArmLSpine CV on the v9 lineage). The key points use SchwingNet, a name,
% not an abbreviation (TITLE.md).
\textbf{Summary statement}: \textbf{\netname{}, trained on rendered CT and radiographs, named lumbar levels on lateral exposures of an unseen C-arm, without a preoperative scan; its transitional naming on real exposures rests on three patients and awaits verified cases.}

\textbf{Key Points}:
\begin{itemize}
\item Never trained on the intraoperative C-arm, \netname{} named the exact vertebra on 595 of 617 detected lateral vertebrae (96.4\,\%; 95\,\% confidence interval: 92.3, 98.4) in 141 cases with cone-beam CT truth, against 276 of 631 anteroposterior vertebrae (43.7\,\%).
\item Counting up from the detected sacrum named 979 of 1,009 bodies on anchored lateral and anteroposterior frames (97.0\,\%); fine-tuned on other cases of that C-arm, \netname{} named 1,079 of 1,159 lateral vertebrae with three in view (93.1\,\%) and 608 of 631 anteroposterior (96.4\,\%), and coupled with the count 991 of 1,009 anchored bodies (98.2\,\%).
\item On the split that chose its checkpoint, its embedding separated the sixth lumbar vertebra from the fifth (silhouette 0.35; 10 persons) more than the fifth from the fourth (0.23); on real exposures its transitional evidence is three cross-validation patients, and certifying 90\,\% exact-level accuracy needs 35 error-free at-risk persons, and seven real-frame candidates exist.
\end{itemize}

\textbf{List of abbreviations}: \printabbreviations

\clearpage
\section*{Abstract}

\textbf{Purpose:} To name vertebral levels from local morphology on one intraoperative C-arm exposure with a network, \netname{}, test it on an unseen C-arm, and count the transitional persons a test needs.

\textbf{Materials and Methods:} Retrospective study of public collections (analysed April--September 2026; intraoperative exposures November--December 2023): 4,093 CT renders, 9,108 annotated radiographs, 4,276 real exposures with volumetric truth; source median ages \RESULT{agesex_csxa_age}--\RESULT{agesex_cobra_age} years, \RESULT{agesex_ts_pct_female}--\RESULT{agesex_csxa_pct_female}\,\% female, none for the intraoperative, cadaveric or phantom sets. A residual encoder and recurrent head named detected bodies, counted from the detected sacrum and coupled. Naming was scored against cone-beam CT in 148 unseen-C-arm cases. Wilson bounds set transitional sample size against pre-stated targets. Intervals were cluster-adjusted; $P$ values, Holm-adjusted exact tests.

\textbf{Results:} In the embedding (checkpoint-selection split), the sixth lumbar vertebra was farther from the fifth (silhouette 0.35; 95\,\% confidence interval: 0.31, 0.39) than the fifth from the fourth (0.23; 0.21, 0.25). Never trained on the C-arm, \netname{} named 595 of 617 detected lateral vertebrae (96.4\,\%; 92.3, 98.4) and 276 of 631 anteroposterior (43.7\,\%; $P<.001$); fine-tuned by case-grouped folds, 601 of 617 (97.4\,\%) and 608 of 631 (96.4\,\%). On anchored frames the sacrum count named 979 of 1,009 bodies (97.0\,\%), the coupled call 991 (98.2\,\%). No real-exposure transitional test person was verified; certifying 90\,\% needs 35 error-free at-risk persons; seven real-frame candidates exist.

\textbf{Conclusion:} \netname{} named ordinary lumbar levels from one lateral exposure of an unseen C-arm; its transitional naming on real exposures rests on three cross-validation patients and awaits verified cases.

\clearpage
% Introduction (CLAIM 2024 items 3 and 4). Every claim about prior work is a row of CLAIMS_LEDGER.md or of
% LEVELCHECK_MODEL.md section 2. No number from this study's own experiments appears here. Last paragraph =
% hypothesis and purpose. Figure 1 (the phenotypes) is cited here first, so the five figures number in the order
% of the story (1 phenotypes, 2 embedding, 3 network, 4 one exposure, 5 persons needed).
% Revision 5 (2026-09-24, REVISION_LOG3.md section 8): the LevelCheck prior-art sentences of paragraph 2 STAY (they
% motivate the problem); the fourth aim of the purpose sentence (the pre-specified comparison) moved to paper 2
% (paper2_levelcheck/MOVED_FROM_PAPER1.md item 7) and is replaced by counting the transitional persons a test needs.
% Revision 2 (2026-09-24, REVISION_LOG2.md): hsiang2011 dropped (DOI not registered; strong2021 carries the sentence);
% Mody "responding", Furuya "consecutive patients" (REFS_AUDIT.md); Huang 2022 worded from its abstract (Long-Film
% images of cadavers; the "about 17 vertebrae" and "tomosynthesis" are not in the abstract); C2 spelled out at first
% use; the network named SchwingNet (TITLE.md).
% Revision 4 (2026-09-24, Greg: thoracic level confirmation; REVISION_LOG3.md section 7): Furuya's regional rates
% (thoracic 0.42 %, lumbar 0.13 %) beside the 0.12 %, and one sentence on why a thoracic level has no
% neurophysiological check (kirshblum2011: no key muscle T2-L1 in the ISNCSCI motor examination; kulik2013: thoracic
% screw stimulation recorded from the rectus abdominis, individual-segment analysis impossible). konin2010 dropped
% with its clause (a review; nakagawa2017 carries the count consequence). Word-level trims for room, no claim changed.

\section*{Introduction}

A vertebral level is not marked; it is counted: fluoroscopy shows three or four bodies, and surgeons count
up from the sacrum or down from the second cervical vertebra (C2) across exposures
\citep{strong2021}. Half of responding spine surgeons report a wrong-level operation \citep{mody2008}, a registry
of 34,544 consecutive patients recorded 0.12\,\% (thoracic 0.42\,\%, lumbar 0.13\,\%) \citep{furuya2026}, and the level first exposed is wrong in 4.3--15\,\% of lumbar
operations before a check film corrects it \citep{dablouk2019,ammerman2006}. No neurophysiological check confirms a thoracic
level: the standard motor examination tests no key muscle between the first thoracic and second lumbar segments
\citep{kirshblum2011}, and thoracic screw stimulation, recorded from the rectus abdominis, cannot separate
segments \citep{kulik2013}. A transitional vertebra --- a
sixth lumbar (L6) or thirteenth thoracic (T13) vertebra, or a fifth lumbar vertebra (L5) fused to the sacrum
(Figure~\ref{fig:phenotypes}) --- changes the count: 8.8\,\% of 3,147 full-length radiographs lacked 24 presacral segments
\citep{sabri2025}, and in 550 lumbar operations the count was wrong in 38 of 71 patients with a lumbosacral
transitional vertebra \citep{nakagawa2017}.

Intraoperative level identification rests on registration: LevelCheck projects a reader's preoperative CT labels
onto the radiograph \citep{levelcheck2012,levelcheck2016}. Its published accuracies are
scored against those labels: 100\,\% label placement \citep{desilva2016}, and 97.6\,\%, five surgeons' median on 62
difficult films with the overlay \citep{manbachi2018}. On a transitional spine the imported
label is only as right as the planning read, and lumbar MRI localisers read without a count from C2 identified the segments in 77.9--82.2\,\% \citep{peh1999}. Registration-free naming on real C-arm images
imported identity from a preoperative radiograph (73.3\,\%) \citep{zhang2023vdvm} or needed long-length (Long-Film)
images of cadavers (88.3\,\%) \citep{huang2022}.

We hypothesised that each transitional vertebra carries local morphology that a network reading a few bodies
can learn, so a level can be named from one C-arm exposure without counting from C2 or importing a label.
The purpose was to build that network (\netname{}), examine what it learns, test it on an unseen C-arm, and
count the transitional persons a test of it needs and the public sources hold.

% Materials and Methods, wrapped around the five figures (rebuilt 2026-09-24 against the figure builds of that
% morning; the sixth figure, the LevelCheck comparison, moved to paper 2 that evening). Every number is traced in
% NUMBERS.md ("The six-figure wrap" and "Revision 2") and sections/methods_sources.md.
% Order here = order of the Results, with the same subheadings (REVIEW2_RSNA.md Fix 25): population and sources,
% phenotypes (Figure 1, Table 2), the reference standard (the automated reference of LEVELCHECK_MODEL.md s14 moved
% to paper 2 in the 2026-09-24 fill pass: no paper-1 result uses it; paper2_levelcheck/MOVED_FROM_PAPER1.md item 8),
% the encoder and its embedding (Figure 2), the network (Figure 3), naming, counting
% and coupling on the intraoperative C-arm (Figure 4) and transitional sample size (Figure 5). First paragraph:
% design, dates, ethics, subject overlap. Last paragraph: statistics. Device-card formation, the corpus manifold,
% the detector's source contributions and detection against angle are in the supplement (Figures S2-S5, Tables
% S1-S4).
% Fill pass (2026-09-24 evening, REVISION_LOG3.md section 8): "Design" dates = April 2026 (month precision: the
% TCIA CT COLONOGRAPHY download scripts are in the repository's first commit, 550cfff, 2026-04-18; no on-grid file of
% that download survives, so the day is not defended) to 24 September 2026. The two CTSpinoPelvic1K readers are
% radiology residents (Greg 2026-09-24 22:15: both PGY2; earlier that day names only, no year was recorded in either repository, so
% none is printed; the descriptor's title page carries their names and department, which the anonymized file may not).
% Revision 5 (2026-09-24, REVISION_LOG3.md section 8; PAPER1_SCOPE.md): the paragraph "The LevelCheck comparison"
% (comparator, strata, unit, s15 detector, s12-s13 secondary analyses), the power sentence of "Statistical
% analysis" and the one-sided .025 clause moved verbatim to paper2_levelcheck/MOVED_FROM_PAPER1.md item 5. What
% stays of that material is paper 1's own cross-validation reporting: the coupled call, the count and the stop
% signal ("The network"; Table 3, Results). "Transitional sample size" is rewritten around the pre-stated Wilson
% targets and defines "at risk" here, since Table 3 no longer does.
% Revision 2 (2026-09-24, manuscript-drive; REVISION_LOG2.md): EVAL_LEAKS.md applied (the zero-shot run is the
% lineage-complete v9 encoder and head, HELD OUT for all 148 cases; the fold models are being re-run from it, their
% numbers PENDING G5); REVIEW2_RSNA.md Fixes 6, 11, 14, 16, 19, 20, 21, 25, 26 and cuts CUT2, CUT5, CUT6, CUT8;
% REVIEW2_STATS.md N1, N4, N12, N13, N17; LEVELCHECK_MODEL.md s14 (automated reference; s15, the primary detector,
% left with the comparison). The automated-reference selection rule is committed (b511985), the rule itself is
% frozen with its SHA-256 before scoring (s14 R2).
% Revision 3 (2026-09-24, collect-and-submit-1; REVISION_LOG3.md): one sentence at the end of "Statistical analysis"
% points to the supplement section "Data integrity and quality controls" (the evaluation gate and every audit tool);
% the abstention rule moved to the naming table's note d for room; small trims measured with check_limits.words().
\section*{Materials and Methods}

\paragraph{Design.}
This retrospective study used de-identified image collections (Table~\ref{tab:data}), analysed
from April 2026, when the first collection was downloaded, to 24 September 2026; intraoperative exposures carry study dates
11 November--12 December 2023. The institutional human research protection office determined this secondary
analysis not to be human participant research (determination 2026-269); consent and the Health Insurance
Portability and Accountability Act do not apply. Every subject was reported by its release or article; CTSpinoPelvic1K also in a
data descriptor \selfcite{ctsp1kdescriptor}.

\paragraph{Sources and transitional phenotypes.}
Rendered CT came from CTSpinoPelvic1K \selfcite{ctsp1k} (801 of 802 records), TotalSegmentator v2 and v3
\citep{totalsegmentator} and VerSe 19/20 \citep{verse2021}; CSXA \citep{csxa}, BUU-LSPINE \citep{buu2023} and
NHANES~II \citep{nhanes} supplied radiographs with released levels or landmarks, and VinDr-SpineXR pseudo-labelled
films \citep{vindr}. Real exposures with three-dimensional truth came from C-ArmLSpine \citep{carmlspine} (296
exposures, 148 of 149 cases, one mobile C-arm, per-case cone-beam CT), DeepFluoro (six specimens
\citep{grupp2020}), COBRA2026 (145 of 467 Elekta cases; seven of 152 downloaded failed a pre-export check
\citep{cobra2026}), a phantom \citep{genevahip} and a cadaveric specimen \citep{ljubljana2d3d}
(Table~\ref{tab:data}; flow, Figure~S1). Transitional persons were typed under a rib
convention as L6, T13, sacralised L5
(Castellvi III or IV \citep{castellvi1984}), Castellvi I/II or lumbar rib (Figure~\ref{fig:phenotypes},
Table~\ref{tab:transitional}): VerSe from its published read; CTSpinoPelvic1K from two radiology residents
(postgraduate year 2) grading independently, six disagreements resolved by consensus
\selfcite{ctsp1kdescriptor}; TotalSegmentator by an independent rib count; intraoperative and COBRA2026 calls are an
automatic labeller's (SPINEPS with VERIDAH \citep{veridah2026}), read by no person.

\paragraph{Reference standard.}
No person read an intraoperative image. Truth is the per-case cone-beam CT, registered to each exposure and
labelled by TotalSegmentator and by SPINEPS with VERIDAH; a vertebra is scored only where both agree (814 of
910), withholding T11 and most of T12, and both count from the sacrum, so agreement is not independent. Naming matches detections one to one to vertebral-body truth boxes
(Table~\ref{tab:naming}); detection is scored against whole-vertebra projections.

\paragraph{The level encoder and what it learns.}
Each body is read from a fixed-scale two-channel crop (image and target marker; target 96 of 320 pixels) by a
ResNet-18 encoder \citep{he2016} into a 512-dimensional embedding. The latest encoder's embedding was mapped by
t-distributed stochastic neighbour embedding (t-SNE) on its validation split after dropping 2,889 crops of 164
persons also in its training keys (Figure~\ref{fig:embedding}); that split chose the checkpoint, so it is not
held out. Separation from the level above is the cosine silhouette, bootstrapped by person and repeated per
render source. The three contested C-ArmLSpine L6 patients, in neither split, and the last free vertebra
of validation persons with an unlabelled fused L5 were placed against the neighbouring clusters (supplement).
Occlusion over ten lateral crops per level from persons outside the training keys gave the share and density
of attribution on the body above, the target and the body below. Renders were styled by image-formation terms
fitted per device, which missed their pre-stated criterion (Figures~S2 and S3, Table~S1).

\paragraph{The network.}
Two detectors never trained on the test C-arm find the bodies (Figure~\ref{fig:network}): a YOLO11m-pose keypoint
detector trained on films boxes the bodies the level model names (confidence 0.35), and a YOLO11m vertebra and
sacrum detector, fine-tuned on real frames of two other devices, boxes the sacrum at confidence 0.74, chosen on
three other devices (Table~S2); 97\,\% detection was pre-stated as counting grade. A two-layer
bidirectional gated recurrent unit (GRU) decodes runs of two to nine embeddings at mixed lengths into a 27-level
posterior. Counting ranks each body from the detected sacrum and reads no morphology. The coupled call chooses,
among the ordinary, L6, T13 and both-variant naming alphabets, the one best fitting the posteriors at the
ordinals; a classifier name differing from the call, or an alphabet margin below the fifth percentile of correct
frames in the other folds, stops the read (Table~\ref{tab:naming}).

\paragraph{Naming, counting and coupling on the intraoperative C-arm.}
C-ArmLSpine was scored zero-shot by the latest encoder and its recurrent head, whose training lists hold no
C-ArmLSpine case; with five case-grouped folds; and, because a case is an
acquisition session rather than a proven person (five dates), with three folds by acquisition day. Per
fold that encoder and head were fine-tuned (15 epochs, learning rate $10^{-4}$), checkpoints chosen on validation
folds excluding every test case (Figure~\ref{fig:exposure}). The primary outcome, designated after a first
zero-shot run was scored, was the exact level of detected lateral vertebrae at full coverage; secondary outcomes
were naming with one or three bodies in view, the count, the coupled call and its stop signal, and detection on
the 56 test frames (Table~S2, Figure~S4). Same-device
detection used nested whole-case subsets of 0 to 210 real frames, and one view alone at 25 and 50 frames (two
seeds; supplement). Detection against view angle used $\pm$90$^\circ$ render pose pools per
device and one patient's 301-frame rotational acquisition, the real pool holding only anteroposterior and lateral
views (Figure~S5). Naming against training persons used NHANES~II, BUU-LSPINE and CSXA with a render-only
encoder on nested person subsets, and C-ArmLSpine per view by case-grouped folds. Render-only detector arms, the fit of accuracy against count, and the
ablation of the rendered rows are in the supplement.

\paragraph{Transitional sample size.}
Pre-stated Wilson targets, a 95\,\% lower bound of 90\,\% for the exact level and 85\,\% for variant sensitivity on
pooled at-risk test persons, set the persons a transitional test needs (Table~\ref{tab:transitional},
Figure~\ref{fig:power}); at risk are the subtypes a count from the sacrum or from C2 misnames: L6, T13 and a
sacralised L5 on a four-lumbar count. The person-disjoint split (50:10:40 per subtype) puts every real transitional
subject on the test side; of its transitional test persons only the three C-ArmLSpine L6 patients, scored in
cross-validation, enter a result here.

\paragraph{Statistical analysis.}
Naming used every available case. Transitional sample size is the smallest person count at which the Wilson
95\,\% lower bound with zero to five errors reaches a target (Figure~\ref{fig:power}a). Intervals are 95\,\% two-thousand-resample bootstraps clustered by case,
patient or person, percentile with the bias-corrected and accelerated (BCa) interval beside; where errors fall in
few clusters, the Wilson interval on the design-effect-adjusted count replaces it. Paired comparisons report the same-item difference with its bootstrap interval and an exact
sign test over clusters, Holm-adjusted over four comparisons named before adjustment (lateral versus
anteroposterior naming; three in the supplement); other contrasts are exploratory; significance was two-sided
$P<.05$. Every reported number passed a patient-keyed evaluation gate
that counts a model's, and its ancestors', training and checkpoint-selection persons as seen (supplement, with
the other integrity checks). Analyses used Python 3.11 and 3.12 (Python Software Foundation,
Wilmington, Del), PyTorch 2.6.0 (PyTorch Foundation), Ultralytics 8.4 (Ultralytics, Frederick, Md), NumPy 2.2.6
and SciPy 1.15.3 and 1.18.1; reporting follows CLAIM 2024 \citep{claim2024}.

% Results, wrapped around the five figures (rebuilt 2026-09-24 against the figure builds of that morning; the
% sixth figure, the LevelCheck comparison, moved to paper 2 that evening). First paragraph = the study population.
% Then the Methods items in order, under the same subheadings (REVIEW2_RSNA.md Fix 25): the phenotypes (Figure 1,
% Table 2), the level encoder (Figure 2), naming, counting and coupling on the intraoperative C-arm (Figure 4,
% Table 3 = tab3_naming.tex; detection in Table S2) and transitional sample size (Figure 5). Every percentage
% carries its denominator here or in the cited table; sources in NUMBERS.md ("Six-figure wrap" and "Revision 2").
% Where the numbers come from (REVISION_LOG2.md, all eval_gate HELD OUT): the zero-shot run is the lineage-complete
% v9 encoder + head (results/eval_gate/leaks/zeroshot_v9.json, overlap 0, all 148 cases); the fold-model, field,
% coupling and stop-signal numbers are the v9-lineage re-run on the PRE truth (results/eval_gate/closure/cv_v9/:
% cv_v9.json, cv_v9_extras_pre.json, coupling_pre/coupling.json; the file Figure 4 draws); the count from the sacrum
% and the detection counts use no level encoder (FIGREC:figure4); the embedding numbers are the v9 selection split
% and say so; the same-device learning curve is results/indomain/INVENTORY.md (HELD OUT rows, body boxes); 673
% agreed AP vertebrae in 143 cases (docs/carmlspine_cv5_mixed_2026-09-21.json#/oracle/AP/kfull), 631 in 142 detected.
% Revision 5 (2026-09-24, REVISION_LOG3.md section 8; PAPER1_SCOPE.md): the paragraph "The LevelCheck comparison"
% and the exact-test sentence of "Transitional sample size" moved verbatim to paper2_levelcheck/MOVED_FROM_PAPER1.md
% item 1; the sample-size paragraph quotes the pre-stated Wilson targets (35 persons at 90 %, 22 carriers at 85 %:
% Table 2 note f, results/power/power.json#n_needed_zero_errors, results/stats_review/recompute.json#R4_wilson). No
% transitional test person of the built split is scored here beyond the three C-ArmLSpine L6 cross-validation cases.
% Fill pass (2026-09-24 evening, REVISION_LOG3.md section 8 "Fill pass"): the study-population sentence quotes
% results/age_sex_by_source.json through \RESULT{agesex_*} (per-source detail in supplement Table S3). The earlier
% draft said the radiotherapy set releases no age or sex; COBRA2026 does (metadata.yaml, ct block: age on 115 and
% sex on 75 of the 145 detector cases), so the sentence now names only the intraoperative, cadaveric and phantom
% sets as releasing none.
% Data-case pass (2026-09-25, REVISION_LOG3.md section 9): the last placeholder (G5) is now the paragraph "Persons,
% films and coverage": film N curve (grid chain 40437811 -> results/coverage/film_ncurve/film_ncurve.json), the
% matched-N coverage arms (chain 40437814 -> results/coverage/experiment/scores.json) and the pooled nested fit
% (tools/indomain_fit_coverage.py -> results/coverage/pooled_fit.json, sensitivity/pooled_fit_conf074.json); every
% number is a \RESULT key of numbers_manifest.json (film_*, cov_*). The sample-size sentences now cite Figure 5a only
% (the old panels c, d moved to paper 2). The case-grouped lateral naming and the sacrum / body-box detection
% numbers of this file are \RESULT keys too (naming_lateral_cv_v9_*, det_carm_*); the rest stay typed until the
% RESULT-conversion pass.
% Acquisition-space pass (2026-09-25 afternoon, REVISION_LOG3.md section 10; Greg: "demonstrate how much they vary in
% the extra dimensions of the fluoro"): Figure 5 is now (a) persons needed, (b) naming vs training persons on film AND
% on the C-arm per view protocol (results/indomain/name_fill.json#figure5b), (c) same-device detection vs real frames
% with the one-view arms (results/indomain/det_fill.json curves.carm + protocol), (d) the matched-N render arms read
% along the lateral angle (experiment/runs/cov_surgio_b51_dr_*/unseen/recall_vs_angle.json). The paragraph "Persons,
% views and angles" walks the panels; the same-device sentences that closed "Naming, counting and coupling" moved into
% it; the pooled fit keeps its numbers without a panel. Keys name_carm_*, det_both_*, det_oneview_*, angle_*.
\section*{Results}

\paragraph{Study population.}
The corpus held 4,093 renders of labelled CT, 9,108 level-annotated films and 4,276 real exposures with
three-dimensional truth, 778 from mobile C-arms (Table~\ref{tab:data}). Of the intraoperative set's 910 cone-beam
CT vertebrae, 814 (89.5\,\%) were agreed, giving 676 agreed lateral vertebrae in 143 cases and 673
anteroposterior in 143. No age, sex or indication is released for the intraoperative, cadaveric or phantom sets.
Across the eight releases that give them (Table~S3), median age ran from \RESULT{agesex_csxa_age} years (CSXA) to
\RESULT{agesex_cobra_age} (COBRA2026) and the share female from \RESULT{agesex_ts_pct_female}\,\%
(TotalSegmentator) to \RESULT{agesex_csxa_pct_female}\,\% (CSXA).

\paragraph{Sources and transitional phenotypes.}
The sources held 194 transitional persons (Table~\ref{tab:transitional}): 110 with L6 (three intraoperative), seven
with T13, 28 with a sacralised L5 (12 with four lumbar vertebrae), 31 with Castellvi I/II, 14 with a lumbar rib and
four unassigned. No released radiograph or C-arm frame carries a human-read T13 or Castellvi grade
(Figure~\ref{fig:phenotypes}).

\paragraph{The level encoder and what it learns.}
On the checkpoint-selection split (13,009 crops of 1,612 persons), the nearest crop from another person was the
same level for 9,740 (74.9\,\%; 95\,\% confidence interval [CI]: 72.2, 77.6) (Figure~\ref{fig:embedding}a, b). L6
separated from L5 (silhouette 0.35; 0.31, 0.39; 98 crops, 10 persons; leave-person-out 0.34) more
than L2--L5 from the level above (0.16--0.23; L5 from L4, 0.23; 0.21, 0.25) and as far as L1 from T12 (0.32; 0.29,
0.35); the classifier named 75 of 98 L6 crops and 111 of 1,012 L5 crops L6 (Figure~\ref{fig:embedding}c; per
render source, supplement). In three
never-seen patients the contested L6 sat nearer the L6 cluster under whole-vertebra boxes only, and in two of four
fused-L5 validation persons the last free vertebra was named L5, the name a count gives (supplement). No T13 crop was outside training. Per unit area, occlusion
attribution was densest on the target body at each of six levels (enrichment 1.39--4.51), but for L4--L6 more fell on
the body below (36--43\,\%) than on the target (24--32\,\%) (Figure~\ref{fig:embedding}d).

\paragraph{Naming, counting and coupling on the intraoperative C-arm.}
The keypoint detector boxed 617 of 676 agreed lateral bodies (91.3\,\%), the held-out detector a sacrum box on 233
of the 286 frames showing one (81.5\,\%) (Figure~\ref{fig:exposure}a). Never trained on C-ArmLSpine, \netname{}
named 595 of 617 detected lateral vertebrae exactly (96.4\,\%; 92.3, 98.4) and 276 of 631 detected
anteroposterior (43.7\,\%), a paired difference of 54.3 points (45.6, 62.9; $P<.001$) on 593 vertebrae of 140 cases
(Table~\ref{tab:naming}). Fine-tuned by case-grouped folds, it named \RESULT{naming_lateral_cv_v9_kn} lateral (\RESULT{naming_lateral_cv_v9_pct}\,\%; \RESULT{naming_lateral_cv_v9_ci}) and
608 of 631 anteroposterior vertebrae (96.4\,\%; 92.5, 98.3), 1.0 point ($-1.9$, 4.0) and 52.6 points (44.1, 61.3)
over zero-shot; by acquisition-day folds, 568 of 617 lateral (92.1\,\%; 87.5, 96.0), 4.4 points below zero-shot
($-9.1$, 0.2). With three lateral bodies in view the case-grouped models named 1,079 of 1,159 (93.1\,\%; 89.2,
96.2), with one 451 of 603 (74.8\,\%) (Figure~\ref{fig:exposure}b); zero-shot, 970 of 1,159 (83.7\,\%) and 385 of 603
(63.8\,\%). On the 1,009 bodies of anchored frames (confidence 0.74), the count
from the sacrum was exact on 979 (97.0\,\%; 94.2, 99.2) (Figure~\ref{fig:exposure}c), the fine-tuned classifier on
986 (97.7\,\%) and the coupled call on 991 (98.2\,\%; 96.5, 99.5), 1.2 points over the count (0.0, 3.2) and 0.5
over the classifier ($-1.5$, 2.4). The read stopped on 23 of 233 anchored frames, catching five of
six erring frames and all four L6 calls: three on five-lumbar frames, one on the four anchored frames of the three
patients one labeller reads six-lumbar (Figure~\ref{fig:exposure}d).
The held-out detector found 231 of 287 vertebrae (80.5\,\%; 74.1, 86.8) and \RESULT{det_carm_sacrum_074_kn} sacra (\RESULT{det_carm_sacrum_074_pct}\,\%), short of the
pre-stated 97\,\% (Table~S2); on the \RESULT{denom_repaired_n} truth boxes that are complete, body-sourced and in the
field (supplement), \RESULT{det_heldout_repaired_kn} (\RESULT{det_heldout_repaired_pct}\,\%; \RESULT{det_heldout_repaired_ci}).
The radiotherapy imager's arc frames carried that detection: alone they found \RESULT{a3_e_only_carm_pct}\,\% of
C-ArmLSpine vertebrae, the cadaver C-arm's frames \RESULT{a3_d_only_carm_pct}\,\%, and adding them raised recall
\RESULT{a3_add_e_to_d_pp} points (\RESULT{a3_add_e_to_d_ci}) (Figure~S4).

\paragraph{Transitional sample size.}
With no error, 35 at-risk persons put the Wilson lower bound at the pre-stated 90\,\% and 22 carriers the
variant-sensitivity bound at 85\,\% (Figure~\ref{fig:power}a). No real-frame transitional test person was
verified: the three C-ArmLSpine L6 patients are scored only in cross-validation (Table~\ref{tab:naming}), and
against 91 unread COBRA2026 test-side screen calls (Table~\ref{tab:transitional}), a two-labeller count with both
junctions in the field found seven at-risk persons with an exported frame, whose error-free Wilson lower bound
would be 64.6\,\% (Figure~\ref{fig:power}a). The split's 51 at-risk CT test persons (most in the current
encoder's training data, Table~\ref{tab:transitional}) would give 93.0\,\% error-free.
\paragraph{Persons, views and angles.}
With a render-only encoder, exact-level naming on films rose with training persons (Figure~\ref{fig:power}b):
NHANES~II \RESULT{film_nhanes_acc_n5_pct} and \RESULT{film_nhanes_acc_all_pct}\,\% at five and all
\RESULT{film_nhanes_n_all_persons} persons; BUU-LSPINE \RESULT{film_buu_acc_n5_pct} and \RESULT{film_buu_acc_all_pct}\,\%
(\RESULT{film_buu_n_all_persons}); CSXA \RESULT{film_csxa_acc_n5_pct} and \RESULT{film_csxa_acc_all_pct}\,\%
(\RESULT{film_csxa_n_all_persons}). Transfer between corpora peaked at
\RESULT{film_transfer_nhanes_to_csxa_pct}\,\% (NHANES~II to CSXA, \RESULT{film_transfer_nhanes_to_csxa_n_films} films).
On the C-arm, naming separated by view rose like film: lateral from \RESULT{name_carm_lat_n0_pct}\,\% with no C-arm
person to \RESULT{name_carm_lat_n50_pct}\,\% with 50, anteroposterior from \RESULT{name_carm_ap_n0_pct}\,\% to
\RESULT{name_carm_ap_n5_pct}\,\% with five and \RESULT{name_carm_ap_nall_pct}\,\% with all; lateral-trained models named \RESULT{name_carm_lat_to_ap_nall_pct}\,\% of
anteroposterior bodies, and training both views cost neither (\RESULT{name_carm_both_lat_pct} and
\RESULT{name_carm_both_ap_pct}\,\%). Same-device detection (body boxes, 0.74) rose from
\RESULT{det_carm_body_recall_074_pct}\,\% (\RESULT{det_carm_body_recall_074_kn}) without the device's frames to a
plateau from 24 frames (lateral \RESULT{det_both_n24_lat_pct}, anteroposterior \RESULT{det_both_n24_ap_pct}\,\%;
\RESULT{det_both_n24_repaired_all_pct}\,\% on the repaired denominator); one-view arms of 25 or 50 frames found
\RESULT{det_oneview_own_min_pct}--\RESULT{det_oneview_own_max_pct}\,\% of its bodies and
\RESULT{det_oneview_other_min_pct}--\RESULT{det_oneview_other_max_pct}\,\% of the other's (Figure~\ref{fig:power}c).
On \RESULT{angle_n_subjects} unseen subjects' renders, render-only detectors at matched count
(\RESULT{cov_arm_surgio_matched_n} renders, \RESULT{cov_arm_surgio_subjects} subjects) trained on anteroposterior
and lateral poses alone found \RESULT{angle_low_trained_pct}\,\% at those angles and
\RESULT{angle_low_untrained_pct}\,\% over the \RESULT{angle_n_untrained_bins} other bins; doubling those frames gave
\RESULT{angle_low2n_trained_pct} and \RESULT{angle_low2n_untrained_pct}\,\%, the same count spread over the arc
\RESULT{angle_high_trained_pct} and \RESULT{angle_high_untrained_pct}\,\% (Figure~\ref{fig:power}d); the Elekta XVI
arms scored zero. On one patient's rotational acquisition the pooled detector found \RESULT{s5_pooled_spin_pct}\,\% of
bodies across the arc, the held-out detector \RESULT{s5_heldout_spin_pct}\,\% (Figure~S5). A nested fit
over \RESULT{cov_pooled_n_points} points preferred count plus lateral-angle coverage at confidence 0.25 (AICc lower by
\RESULT{cov_pooled_daicc}; $P\RESULT{cov_pooled_lr_p}$), not at 0.74 ($P\RESULT{cov_pooled_lr_p_conf074}$).


\section*{Discussion}
% Discussion, first part: the problem and the primary results, what the encoder learned, what the real-frame
% transitional evidence here is (and that the powered comparison is a separate study), and relation to other
% studies. Every comparator is a row of CLAIMS_LEDGER.md; no table or figure is cited in the Discussion; no claim of
% priority (journal rule; REVIEW_RSNA.md Fix 6).
% Revision 5 (2026-09-24, REVISION_LOG3.md section 8): paragraph 3 (LevelCheck's published accuracies, the frozen
% comparator, the 152 persons) moved verbatim to paper2_levelcheck/MOVED_FROM_PAPER1.md item 6 and is replaced by
% the plain statement of what the transitional evidence on real exposures rests on; paragraph 4's "comparison arm"
% is named for what it is (a level model trained from ImageNet weights on the split's training persons).
% Revision 2 (2026-09-24, REVISION_LOG2.md): only held-out numbers here (EVAL_LEAKS.md): the zero-shot v9 run
% (96.4 %) and the count from the sacrum (97.0 %, no level encoder); no fine-tuned number is credited to the network
% that never saw the device (DRIVE_CHECK D2; REVIEW2_RSNA Fix 3; REVIEW2_STATS N2), and the fine-tuned fold results
% are PENDING G5, so none is quoted here (the journal: no results of incomplete investigations in the Discussion).
% REVIEW2_RSNA S1 (paragraph 2), CUT1 (paragraph 3), A1d (paragraph 4, without its fine-tuned numbers).
% Revision 3 (2026-09-24, collect-and-submit-1; REVISION_LOG3.md): two clause-level trims in paragraph 4 for room
% ("the render arc pools span it"; "paired with its cone-beam CT" dropped); no number or claim changed.
% Revision 4 (2026-09-24, REVISION_LOG3.md section 7): the diffpose2024 future-work clause at the end of paragraph 4
% dropped (its slot and ten words went to the Introduction's thoracic sentence); zheng2018pda's number unchanged.
% Film N-curve ablation (2026-09-26, REVISION_LOG3.md section 13; jobs 40451157 imagenet / 40451158 filmonly): the
% clause "A level model trained from ImageNet weights saturated likewise in development (supplement)" is replaced by
% the measured ablation, which IS that comparison done properly - the ImageNet arm on the same merged bank, plus a
% film-only arm with every render row removed. Both directions are stated: renders help where film is scarce
% (7 of 9 pre-stated cells separated on every seed) and one corpus ends 1.6 points the other way once every person is
% used. All five numbers are \RESULT{} keys into results/coverage/film_ncurve/ablation/ablation_read.json; none typed.
% Room came from the last sentence of Methods "Naming, counting and coupling", compressed to its supplement pointer.
% Acquisition-space pass (2026-09-25 afternoon, REVISION_LOG3.md section 10): paragraph 4 now states what Figure 5b-d
% show (the plateau from 24 frames is the same pooled or separated by view; frames of one view buy that view; on
% renders accuracy sits where the training poses sit) and the reading "accuracy followed where the frames sat in
% acquisition space, not how many there were"; the pooled fit's clause kept, its "held-out-angle effect" sentence
% replaced by the one-view and per-angle evidence. No figure cited (journal rule).

Levels are counted across partial exposures, and a transitional vertebra makes the count and an imported label
disagree; \netname{} names a level from one C-arm exposure by the shapes of the bodies in view. Never trained on the intraoperative C-arm, it named 96.4\,\% of
detected lateral vertebrae exactly, and counting from the detected sacrum was exact on 97.0\,\% of anchored
bodies; the anteroposterior view did not survive the unseen device until fine-tuned on it (96.4\,\%); the lateral
did not measurably change.

The encoder set a sixth lumbar vertebra apart from the fifth as far as the rib boundary sets L1 apart from T12,
and per unit area drew most on the target body, as local identity predicts. The evidence is thin: ten persons on a
selection split, an order reversed in one render source, a contested L6 in three never-seen patients holding
under one box convention only, and in two of four fused-L5 persons the name a count gives.

On real exposures the transitional evidence here rests on three C-ArmLSpine patients whose sixth lumbar vertebra
one labeller reads and the other does not, scored in cross-validation and agreeing with that read on one of five
bodies; that is a description, not a test of transitional naming. The powered comparison with registration carrying
a planning read, on verified at-risk persons against a comparator fixed before scoring, is a separate study. No
accuracy here is a wrong-level surgery rate.

Detection, not naming, lost most to the unseen device; same-device frames restored it by 24 frames to a plateau
set by truth boxes cut at the edge of the field, not by the detector. Separating those frames by view neither raised
nor moved the plateau; frames of one view bought that view alone. On renders the same held along the angle:
detectors trained at two poses found bodies there and not between, and the same count spread over the arc held at
every angle. Accuracy followed where the frames sat in acquisition space, not how
many there were. Rendered CT made a few films enough: at five to twenty-five persons per corpus the render-pretrained
encoder named ordinary levels \RESULT{abl_gap_min_pts} to \RESULT{abl_gap_max_pts} points above film alone, separated
on every seed in \RESULT{abl_cells_separated} of \RESULT{abl_cells_total} cells; with every person the arms met, one
corpus \RESULT{abl_worst_delta_abs_pts} points the other way (supplement). Styled renders stayed separable from real frames (supplement); paired domain adaptation
\citep{zheng2018pda} is untried here.

% the blank line below is load-bearing: discussion_limits.tex opens with a comment line, which swallows the
% newline, so without it the limitations paragraph would run on from the last Discussion paragraph in the PDF

% Discussion, second part: the limitations (second-to-last paragraph) and the summary (last paragraph).
% Revision 2 (2026-09-24, REVISION_LOG2.md): the deleted training record of the first encoder (EVAL_LEAKS.md U1) is a
% limitation; REVIEW2_RSNA C10 and REVIEW2_STATS N11 ("at least 19"; "untested"); the embedding claim carries its
% selection split. REVIEW2_RSNA B4 (day-grouped fine-tuning below zero-shot) waits for the G5 re-run: add it here if
% the v9-lineage day-grouped result again falls below the zero-shot 96.4 %.
% Revision 6 (2026-09-24 ~22:15, Greg: "this limitation is unacceptable"): the sentence on the first encoder's deleted
% training record removed. Every number the paper reports comes from the lineage-verified encoder (Table 3 note g), so
% the paragraph no longer describes runs the paper does not report; the audit trail stays in REVISION_LOG2/3,
% NUMBERS.md and results/eval_gate/. The day-grouped sentence stays.
% Revision 5 (2026-09-24, REVISION_LOG3.md section 8): the two comparison clauses ("so the comparison uses a
% separately trained arm"; the summary's "better than registration ... at least 19 verified at-risk persons") moved
% to paper2_levelcheck/MOVED_FROM_PAPER1.md item 6; the summary now states the real-frame transitional evidence.
This study has limitations. One intraoperative device is held, and its truth is two automatic labellers
agreeing, which removes the thoracolumbar junction and cannot settle a transitional count; whether its 148 cases
are 148 people cannot be established, so every intraoperative interval, treating cases as independent, is too
narrow if they are not. Day-grouped fine-tuning scored below zero-shot laterally.
The embedding is of a split that chose the checkpoint; most
transitional test persons of the built split were in the current encoder's lineage, so no model reported here scores them.
Held-out detector numbers use whole-vertebra boxes; the posterior is uncalibrated; nothing has been evaluated prospectively.

In summary, \netname{} named ordinary lumbar levels on lateral exposures of an unseen C-arm, and on its selection
split its embedding set a sixth lumbar vertebra apart from the fifth. Its transitional naming on real exposures
is unverified: three cross-validation patients carry the only real transitional exposures scored, certifying
90\,\% needs 35 error-free at-risk persons, and seven candidates exist. The clinical role is decision support abstaining visibly beside the surgeon's
count, never replacing the check film.


\section*{Other information}
% Revision 2: REVIEW2_RSNA.md Fix 2 (VinDr-SpineXR is under a PhysioNet data use agreement; the literature panels are
% figures of published articles, not a release), Fix 1 (the case134 directory is used for Figure S5, not excluded),
% Fix 6 (Geneva licence); "JSON" left the text (DRIVE_CHECK item 7); "the author" left the writing-tools line
% (REVIEW2_RSNA.md Fix 5: no statement of single authorship in the anonymized file).
\textbf{Data availability.} Every image source is listed in Table~\ref{tab:data} with its licence and download source (flow of cases in Figure~S1 of the supplement); VinDr-SpineXR requires a PhysioNet data use agreement, and the literature panels are not redistributed. The derived frame layout is not redistributed because the licences differ; it is rebuilt from the sources by the released exporters. The fitted image-formation parameters are released as one file per device. Result files behind every number, with the path of each number inside them, are listed in the repository's claims ledger and in the file NUMBERS.md.

\textbf{Code and model availability.} Code, the fitted device parameters and the trained weights are in a private repository (branch L6\_Class) available to the reviewers and editors on request; a public deposit with a revision identifier and a registered weight hash will accompany publication. The weights are released for retrospective analysis and for bench, phantom and cadaver evaluation; nothing here has been evaluated prospectively.

\textbf{Ethics.} Secondary analysis of de-identified imaging, determined by the institutional human research protection office not to be human participant research (determination 2026-269). C-ArmLSpine and COBRA2026 are used under their own releases (CC BY 4.0 and CC BY-NC 4.0), the Geneva phantom under CC BY-SA 4.0, and the DeepFluoro and Ljubljana cadaveric sets under CC BY-NC-SA 4.0 and CC BY-NC-ND. VinDr-SpineXR was obtained under its PhysioNet data use agreement and is covered by a separate determination of the same office (\PENDING{G0}{VinDr determination number and date}); it trained the intraoperative fold models only (Table~\ref{tab:data}). The literature fluoroscopy panels are figures of published articles, used for training only and not redistributed.
% Fill pass (2026-09-24 evening): the HPR memo 2026-269 (2 September 2026) scopes itself to data downloadable
% "without an application, request for permission, eligibility review, or data use agreement" and says a
% DUA-governed dataset "falls outside this determination and a separate IRB submission will be made before that
% dataset is used"; VinDr landed 2026-09-07. 2026-09-25 (Greg): the separate determination is being filed and is to
% be assumed granted, so the sentence says covered and the one \PENDING token left is its number and date (the
% "or its removal from the fold banks" wording is gone). The clause "and the dataset authors were notified on
% <date>" stays dropped: no record exists that a notice was sent (finding dated 2026-09-22; Zenodo record
% 10.5281/zenodo.22768651 is the contact); if Greg sends it, the clause returns with the date.

\textbf{Registration.} Not a clinical trial; not registered.

\textbf{Funding and acknowledgments.} Stated on the full title page, which is submitted separately for double-anonymized review.

% Guarantor and conflict-of-interest statements removed from the anonymized file (Packet phase, 2026-09-24): both
% are on the full title page (packet_src/build/full_title_page.*), and the cover letter says they are not here
% (REVIEW2_RSNA.md blocking item 5; packet_src/statements.md EDIT_PHASE_NOTES 2).

\textbf{Use of writing tools.} Drafting, figure and table assembly and the number-tracing checks were assisted by a large language model operating on the result files; responsibility for the content rests with the authorship stated on the full title page.

\rsnabibliography

\clearpage
\section*{Figure Legend(s)}
\makeatletter\@starttoc{lof}\makeatother

% Figure 1 (SchwingNet/tools/figure1.py; panels tools/figure1_panels.py on the grid; record figures/figure1.record.json).
% Revision 2: the legend adds age and sex (REVIEW2_RSNA Fix 18; figure1.record.json components[].demographics and
% film_demographics) and so differs from figures/figure1.caption.txt (08:24), which carries none; the figure owner
% should adopt this legend. (c) keeps its own L5 label, (d) carries the fused body in the sacrum label (the at-risk
% distinction of Table 2 note f).
\begin{figure}[tbh!]
\centering
\RESULTFIG{figure1}
\caption{Transitional phenotypes. (a) Sixth lumbar, 54-year-old man; (b) thirteenth thoracic, 24-year-old woman; (c) sacralised L5, own label, Castellvi IV, 89-year-old, sex other; (d) four lumbar, L5 in sacrum, Castellvi IIIb, 60-year-old man. Rows: labelled CT (sagittal; b, coronal slab; c, surface); render; another patient's film (women, 31, 64). Arrow and red, transitional vertebra; bar, 20 mm; hatched, none human-read.\label{fig:phenotypes}}
\end{figure}

% Figure 2 (SchwingNet/tools/figure2.py + figure2_analyse.py; grid tools/fig2_heldout_embed.py, fig2_occlusion_clean.py;
% record figures/figure2.record.json). The validation split is the encoder's selection split: not held out.
\begin{figure}[tbh!]
\centering
\RESULTFIG{figure2}
\caption{Level-encoder embedding, validation crops (checkpoint-selection split). (a) t-SNE (t-distributed stochastic neighbour embedding) by level; triangles, sixth lumbar vertebra (L6); diamonds, last free vertebra above fused L5; stars, one labeller's L6 in never-seen patients; rings, their other vertebrae; T13, thirteenth thoracic. (b) Lumbosacral window. (c) Silhouette; CI, confidence interval. (d) Occlusion enrichment: attribution over area share; ten lateral crops per level.\label{fig:embedding}}
\end{figure}

% Figure 3 (SchwingNet/tools/figure3_network.py; record figures/figure3.record.json). A schematic: no result on it.
\begin{figure}[tbh!]
\centering
\RESULTFIG{figure3}
\caption{\netname{}; shaded blocks learned, sizes in pixels. (a) Two detectors, neither naming a level; one boxes bodies, one the sacrum. (b) Level model: two-channel crop per body, ResNet-18, bidirectional gated recurrent unit (GRU) across bodies, 27-level posterior; red, T13 (thirteenth thoracic) and L6 (sixth lumbar). (c) Ordinals counted from the detected sacrum. (d) Posteriors choose the alphabet; outlined, the call.\label{fig:network}}
\end{figure}

% Figure 4 (SchwingNet/tools/figure4_frame.py; record figures/figure4.record.json, rebuilt 2026-09-24 ~10:41 from
% the v9-lineage case-grouped folds on the PRE truth, every drawn input eval_gate HELD OUT; models "v9init_pre").
% Revision 2 round 2: legend = figures/figure4.caption.txt of that build, verbatim (60 words); the Results quote the
% same closure file (results/eval_gate/closure/cv_v9/cv_v9.json, coupling_pre/coupling.json).
\begin{figure}[tbh!]
\centering
\RESULTFIG{figure4}
\caption{Reading C-arm exposures. (a) Bodies (blue) and sacrum (dashed) boxed by detectors never trained on this C-arm. (b) Field b, enlarged, named by shape, fine-tuned on other cases (``none'': no name). (c) Ordinals from the sacrum, ordinary alphabet. (d) A frame one CT labeller reads six-lumbar: shape fits the L6 alphabet (outlined); a stop, not a verdict. L1--L6, lumbar vertebrae.\label{fig:exposure}}
\end{figure}

% Figure 5 (SchwingNet/tools/fig_power.py -> figures/figure5.{pdf,png,tif}, figure5.caption.txt, figure5.record.json,
% results/power/power.json). Rebuilt 2026-09-25 afternoon (REVISION_LOG3.md section 10) around what accuracy follows
% in acquisition space: (a) the Wilson persons-needed chart with the two pre-stated targets (90 % exact level, 85 %
% variant sensitivity); (b) naming against training persons on the three film corpora (film_ncurve.json) AND on the
% C-arm per view protocol (results/indomain/name_fill.json#figure5b: own view solid, lateral-trained read on
% anteroposterior dashed, both views x); (c) same-device detection against real frames of the device
% (results/indomain/det_fill.json: both-view arm per test view; one-view arms of 25 and 50 frames as triangles, filled
% on their own view, open on the other); (d) the matched-N render-only arms read along the lateral angle on 40 unseen
% subjects (experiment/runs/cov_surgio_b51_dr_*/unseen/recall_vs_angle.json). The 10:26 panels (c) arm dot plot and
% (d) pooled-fit scatter are gone; the pooled fit's numbers stay in the Results paragraph. Every figure is included
% through \RESULTFIG{figureN} (numbers_manifest.json "figures"; a red pending marker if the file is missing). The
% legend below is figures/figure5.caption.txt verbatim, LaTeX-escaped (58 words).
\begin{figure}[tbh!]
\centering
\RESULTFIG{figure5}
\caption{Persons, views, angles. (a) Wilson 95\,\% lower bound by errors against transitional test persons; lines, targets. (b) Naming against training persons: film corpora and one C-arm by view (dashed, lateral-trained read on anteroposterior; x, both views). (c) Same-device detection against real frames; triangles, one-view arms (open, other view). (d) Render-only detectors at matched N against lateral angle.\label{fig:power}}
\end{figure}

% ---- Tables 1-3, in citation order (the journal: tables after the legends, one per page, numbered in order of
% first citation). 1 = image sources (tab:data), 2 = transitional persons by subtype (tab:transitional),
% 3 = naming, counting and coupling (tab:naming; file tab3_naming.tex). The former Table 3, the frozen LevelCheck
% comparator (tab_levelcheck.tex), moved to paper2_levelcheck/tables/ on 2026-09-24. Each .tex has a .md twin and a
% .sources.md naming the file behind every cell; fit check slurm/tables_check.sh + tools/tables_check.py
% (tables/tables_check.tex inputs the same three).
% Table 1 -- every collection the models learned from or were tested on, one row per collection; no merged cells
% (group headings sit in the first cell of their own row). Every cell is traced in tab_corpus.sources.md; the same rows
% are in tab_corpus.md. Re-derived 2026-09-23 against docs/DATA_COMPLETENESS_AUDIT.md and docs/DATASHEET.md sections
% 2.1, 2.6 and 3: NHANES II is 543 annotated films of a 17,075-entry release; COBRA2026 is 145 patients used of 692
% public cases (467 Elekta + 225 Varian), the rest fetching (G0); C-ArmLSpine ships 299 spot exposures and two
% cone-beam sweeps, of which 296 frames of 148 cases were scored (case134's two exposures missed by a folder walker,
% case118's oblique scored as a lateral -- G0 repair); VinDr's three counts are three objects and 2,071 of its films
% sit as 2,632 crops in the intraoperative fold banks (none in the film bank); FluoroCap likewise (950 panels as
% 2,100 crops). New column "Released from it" (corpus / parameters / weights; footnote k).
% Revision 2 (2026-09-24, REVISION_LOG2.md): title "one row per collection" (REVIEW2_RSNA Fix 2: VinDr-SpineXR is under
% a data use agreement, FluoroCap is figures of articles); gate labels and pipeline slang out of the notes (Fix 26);
% footnote c widened to the implant-render key alias (closure MANUSCRIPT_CORRECTIONS.md 2.4); Geneva and Ljubljana R2
% are within-subject and footnote k says what each fold leaves out (closure 3b; EVAL_LEAKS L6); the film rows'
% validation split named as the checkpoint-selection split (closure 3c).
% Layout: landscape. The whole table (caption, tabular, notes) sits in a 23.2 cm minipage rotated 90 degrees with
% graphicx's \rotatebox inside the float, so no package is added to rsnastyle.tex and endfloat still moves it to the
% end. The rotated box is 23.2 cm tall (text height 23.7 cm) and its width is the table's natural height, which must
% stay under the 16 cm line: slurm/tables_check.sh + tools/tables_check.py measure that and every font size (7 pt
% floor at native size, the size a landscape page prints it). Footnote markers are \textsuperscript so a Word
% conversion keeps them as superscripts. Body 7 pt on 7.6 pt.
\begin{table}[p]
\centering
\begingroup
\singlespacing
\fontsize{7}{7.6}\selectfont
\newcommand{\fnm}[1]{\textsuperscript{\fontsize{7}{7}\selectfont #1}}
\setlength{\tabcolsep}{2pt}
\renewcommand{\arraystretch}{1.0}
\rotatebox[origin=c]{90}{\begin{minipage}{23.2cm}
\captionsetup{font=footnotesize,skip=4pt}
\centering
\caption{Image sources: one row per collection, with its modality and device, its role, the unit the split keeps disjoint, what was scored, its truth and what this work releases from it.}
\label{tab:data}
% column widths sum to 22.10 cm; with 14 gaps of 2 pt (0.99 cm) the tabular is 23.09 cm in a 23.1 cm minipage
\begin{tabular}{%
>{\raggedright\arraybackslash}p{2.20cm}%
>{\raggedright\arraybackslash}p{2.50cm}%
>{\raggedright\arraybackslash}p{3.20cm}%
>{\raggedright\arraybackslash}p{3.30cm}%
>{\raggedright\arraybackslash}p{3.70cm}%
>{\raggedright\arraybackslash}p{3.40cm}%
>{\raggedright\arraybackslash}p{3.80cm}}
\toprule
Collection\fnm{j} & Modality; device & Role in this work & Distinct subjects (split unit) & Frames, films or renders (train / val / test) & Truth source & Released from it\fnm{k} \\
\midrule
\emph{CT with labels, rendered} & & & & & & \\
CTSpino\-Pelvic1K & abdominopelvic CT & level bank\fnm{a} & 801 patients of 802 records\fnm{b} (TCIA patient identifier) & 801 renders & release labels: levels, ribs, pelvis & corpus; level weights \\
Total\-Segmentator v2 and v3 & whole-body CT & level bank; styled corpus (142 subjects, 17 validation) & 1,935 scans\fnm{c} (CT scan) & 1,170 (v2) and 1,735 (v3) renders & labels mapped by name & corpus; level weights \\
VerSe 19/20 & spine CT & level bank; no accuracy quoted\fnm{c} & 355 subjects, 374 series (subject) & 387 renders; 374 soft-tissue renders & VerSe masks C1--L6; published transitional table & corpus; level weights \\
\emph{Radiographs} & & & & & & \\
CSXA & cervical lateral & level bank; checkpoint selection & 4,963 individuals (individual) & 4,963 films & per-vertebra levels, released & corpus; level weights \\
BUU-LSPINE & lumbar AP and lateral & level bank; checkpoint selection & 3,600 cases\fnm{d} (case number) & 3,603 lateral films; the 3,600 AP films in no bank & end-plate lines L1--S1; levels by position & corpus; level weights \\
NHANES II & lateral survey film & level bank; checkpoint selection & 364 persons among the annotated (sample person) & 543 annotated films (542 on disk) of 17,075 release entries\fnm{e} & landmarks, released & corpus; level weights \\
VinDr-SpineXR & AP and lateral & intraoperative fold banks only, training, pseudo-labelled\fnm{f} & 5,000 studies\fnm{c} (study) & 2,071 films as 2,632 crops, of 10,466 released images & lesion boxes released; levels pseudo-labelled & corpus (all 10,466 images) \\
FluoroCap & literature fluoroscopy panels & intraoperative fold banks only, training\fnm{f} & not established\fnm{g} (source article or video) & 950 panels as 2,100 crops, of 29,484 harvested & none released & none: dropped \\
\emph{Real exposures with three-dimensional truth} & & & & & & \\
C-ArmLSpine & mobile C-arm, Surgio-B51; per-case CBCT & naming cross-validation; held-out detector test & not established; 149 released cases, 148 scored\fnm{h} (case = acquisition session) & 296 scored: 210 / 30 / 56 (105 / 15 / 28 cases); released 299 spot exposures, two cone-beam sweeps\fnm{h} & CBCT, two automatic labellers, agreed levels; pose by registration & corpus; level weights (fold models); no detector weight \\
DeepFluoro & mobile C-arm, Cios Fusion; cadaver CT & detector training; device parameters & 6 specimens (specimen; 4 / 1 / 1) & 366 frames: 238 / 24 / 104 & CT and C-arm pose per frame & corpus; parameters (40 of 366 frames, R² 0.905 vs 0.899); detector weights \\
COBRA2026, Elekta XVI & radiotherapy kV imager; planning CT & detector training; device parameters & 145 patients of 467 public Elekta cases\fnm{i} (case = patient; 110 / 8 / 27) & 3,480 frames: 2,640 / 192 / 648 & planning CT, deformably registered; RTK geometry & corpus (all 692 public cases); parameters (40 of 335, R² 0.974 vs 0.964); detector weights \\
Geneva hip & mobile C-arm, BV Pulsera; phantom CT & unseen device, test only; device parameters & 1 phantom (none: one object) & 116 frames, all test & CT and gold-standard pose & corpus; parameters (40 of 116, within-subject R² 0.952 vs 0.917) \\
Ljubljana 2D-3D-GS & turntable flat panel, PIXIUM 4600; cadaver CT & device parameters only & 1 specimen (not split) & 18 frames & CT and gold-standard pose & parameters (14 of 18, within-subject R² 0.679 vs 0.650); no frame \\
\bottomrule
\end{tabular}
\par\vspace{1pt}
\begingroup\fontsize{7}{7.6}\selectfont\raggedright
AP = anteroposterior; CBCT = cone-beam CT; CDC = Centers for Disease Control and Prevention; RTK = reconstruction toolkit; TCIA = The Cancer Imaging Archive; level bank = the training set of the level encoder and head; R² = coefficient of determination.
\fnm{a}Physics-styled re-renders of the three CT collections (7,045 over 3,029 subjects; 1,482 over 81 transitional subjects) repeat subjects already counted; bank subject counts are not additive.
\fnm{b}Record 0533 was a truncated download and is in no trained bank; rebuilt and replaced in the release (802 records in the released corpus, supplement Table~S3), it enters no reported number.
\fnm{c}No patient identifier in the release: the scan, series or study is the finest recoverable unit; the TotalSegmentator count folds 24 renumbered pairs and three duplicate groups; in the level banks, CTSpinoPelvic1K implant renders were keyed apart from their source CT and, in the first bank, VerSe series as subjects, which put CT-derived persons on both sides of each bank's split (192 in the first, 164 in the second; none from a film corpus), so no CT-render naming accuracy is quoted.
\fnm{d}3,598 by content: two case pairs share byte-identical films, which the built split merges (supplement Table~S3).
\fnm{e}The CDC index lists 17,075 films; two are served empty at the source, so 17,073 are on disk; one annotated film survives only inside the annotation archive.
\fnm{f}Fold banks = the capped per-fold crop banks of the intraoperative cross-validation (about 31,400 crops per fold); the film bank behind the plain-film accuracies holds no VinDr or FluoroCap image. The released corpus takes all 10,466 VinDr images and drops FluoroCap (a republished figure is the same patient).
\fnm{g}A panel keys to its source article or video, which may hold several patients or one patient republished.
\fnm{h}The archive ships 299 spot exposures over 149 cases (149 AP, 149 lateral with one repeat, one oblique) and two 301-projection cone-beam sweeps; the scored bank holds one AP and one lateral of 148 cases (case134's two, filed in its cone-beam directory, were missed by the export script) and scores case118's −58° oblique as a lateral. No patient tag survives; the cases fall on five acquisition dates.
\fnm{i}692 public cases: 467 Elekta XVI (centres A--D), 225 Varian OBI (F--G); 152 Elekta cases were fetched before the analysis and 145 kept --- two pilot cases (A004, A005) never entered the export list and five (A081, A239, A260, A303, B000) failed an axis-map check, B000 because its planning CT is released outside the imaged volume; 26 Varian cases trained no model; the remaining cases enter no reported number.
\fnm{j}Totals over the additive rows: 4,093 renders (801 + 1,170 + 1,735 + 387); 9,108 level-bank films (4,963 + 3,603 + 542; 2,071 VinDr films and 950 FluoroCap panels sit in the fold banks alone); 4,276 real frames scored (296 + 366 + 3,480 + 116 + 18), of which the mobile C-arm rows (296, 366, 116) give 778; subjects are not totalled because the units differ.
\fnm{k}Corpus: rebuilt from the public release by the released exporters (the licences differ, so derived frames are not redistributed). Parameters: a fitted nine-term image-formation file for the device, with the frames it was fitted on and its median R² over six folds on the frames each fold left out, against uncalibrated physics (patients left out for Cios Fusion and Elekta XVI; views or frames of the one object for BV Pulsera and PIXIUM 4600). Weights: the released weights were trained on it (level weights = the fold models of the intraoperative cross-validation; detector = the held-out-device detector trained on COBRA2026 and DeepFluoro real frames); corpus and weights are released at publication. Downloaded but used by no trained component: SpinalAI, AASCE, SpineSAC phantom frames, Mendeley-716 and one internal render set.
\par\endgroup
\end{minipage}}
\endgroup
\end{table}

% Table 2 (file name kept from the four-table layout; it prints second because Methods cites it second) -- the
% transitional persons by subtype: what each subtype is, how many persons carry it per source, what truth the subtype
% rests on, where the built split puts them, how many of them the current encoder trained on, the unread COBRA2026
% screen calls, and the best-case Wilson bound their test count allows. Rewritten 2026-09-23 around Figures 1 and 5:
% the claim columns of the earlier table (sensitivity, specificity, exact level, registration ceiling, simulated
% surgeon) are gone; the comparison (formerly Table 3, tab_levelcheck.tex, and Figure 6) is paper 2 since 2026-09-24
% (paper2_levelcheck/). No merged cells. Sources, one line per cell, in tab4_transitional.sources.md; the same rows
% in tab4_transitional.md.
%   results/rare_class/transitional_subtypes.csv  persons per subtype and source (distinct patient_key; 196 rows, 194 persons)
%   results/rare_class/split_plan.md               the built G3 split (make_split.py apply 2026-09-23 20:18, seed 0,
%                                                  leak audit PASS): test persons per subtype, and val+test persons in
%                                                  the v8 or v9 level-bank train split. OWNED BY THE SPLIT WORKFLOW.
%   results/power/inputs/cobra_lstv_screen.grid.json  the COBRA2026 screen calls (SPINEPS+VERIDAH, unread), as Figure 5 draws them
%   figures/figure5.record.json (panel d)          Wilson lower bounds at zero errors; recomputed here (z = 1.959964)
% Updated 2026-09-24 to the rebuilt Figure 5 (figures/figure5.record.json): the COBRA2026 column adds, in brackets,
% the calls their own field supports (57 of 91: L6 11, T13 38, four-lumbar 8), the Wilson column their best case,
% and footnote c said the three real C-ArmLSpine L6 are excluded from the comparison's primary analysis in advance
% (LEVELCHECK_MODEL.md s4), so their 43.9 % bound is no longer printed as an available real-frame bound.
% Revision 5 (2026-09-24, REVISION_LOG3.md section 8): the comparator clauses of notes c, d and f moved to paper 2
% (paper2_levelcheck/MOVED_FROM_PAPER1.md item 3); note f now defines "at risk" itself; note c points to the
% cross-validation row of the naming table.
% Revision 2 (2026-09-24, REVISION_LOG2.md): the COBRA2026 bracket and the "real" Wilson bounds follow the rebuilt
% Figure 5 (figures/figure5.record.json, 08:26; results/transitional_registry/cobra/rows.csv built 06:49): persons a
% two-labeller count (SPINEPS+VERIDAH and a TotalSegmentator rib count, both junctions in the field) calls
% transitional AND that have an exported frame (L6 3, four lumbar 4, T13 0; pooled at risk 7), replacing the 57
% "supported" screen calls; the Any row is the frozen at-risk strata (L6, T13, four-lumbar sacralised: 51 CT test
% persons, Wilson 93.0) with the own-label sacralised L5 as naming unaffected (REVIEW2_STATS N7); footnote e adds the
% labeller's T13 record, 0 of 52 on CTSpinoPelvic1K (results/spineps_ids/IMPACT.md line 230); footnote d drops the
% bank names (REVIEW2_RSNA Fix 26).
% Layout: landscape (rotated 23.2 cm minipage, see tab_corpus.tex); 7 pt on 7.6 pt.
\begin{table}[p]
\centering
\begingroup
\singlespacing
\fontsize{7}{7.6}\selectfont
\newcommand{\fnm}[1]{\textsuperscript{\fontsize{7}{7}\selectfont #1}}
\renewcommand{\arraystretch}{1.0}
\rotatebox[origin=c]{90}{\begin{minipage}{23.2cm}
\captionsetup{font=footnotesize,skip=4pt}
\centering
\caption{Transitional persons by subtype: definition, persons per source, the truth behind the subtype, the test side of the built person-disjoint split, the share of held-out persons the current level encoder trained on, the unread COBRA2026 screen calls, and the best-case Wilson 95\,\% lower bound the test count allows.}
\label{tab:transitional}
% widths sum to 21.90 cm; 14 gaps of 1.9 pt (0.94 cm) -> 22.84 cm
\setlength{\tabcolsep}{1.9pt}
\begin{tabular}{%
>{\raggedright\arraybackslash}p{2.30cm}%
>{\raggedright\arraybackslash}p{4.50cm}%
>{\raggedright\arraybackslash}p{2.70cm}%
>{\raggedright\arraybackslash}p{3.80cm}%
>{\raggedright\arraybackslash}p{2.10cm}%
>{\raggedright\arraybackslash}p{2.30cm}%
>{\raggedright\arraybackslash}p{1.90cm}%
>{\raggedright\arraybackslash}p{2.30cm}}
\toprule
Subtype\fnm{a} & Definition; what a detector can see\fnm{b} & Persons: total (CTSP / VerSe / TS / C-Arm) & Truth behind the subtype & Test persons, built split\fnm{c}: CT + real & Validation and test persons the current encoder trained on\fnm{d} & Unread COBRA2026 screen calls (two-labeller count, with a frame)\fnm{e} & Wilson lower bound with no error at the test count, \%\fnm{f} \\
\midrule
L6 & six rib-free vertebrae above the sacrum; six lumbar bodies where the sixth carries a label & 110 (18 / 86 / 3 / 3) & CTSP two-reader consensus or release label; VerSe published read (38 persons on the audit read alone); TS independent rib count; C-Arm one automatic labeller & 43 + 3 & 49 of 54 & 38 (3) & 91.8 (43 CT); 43.9 (3 real) \\
\addlinespace[2pt]
T13 & thirteen rib-bearing vertebrae & 7 (0 / 6 / 1 / 0) & VerSe mask label 28; TS independent rib count & 3 + 0 & 4 of 4 & 45 (0) & 43.9 (3 CT); none (0 real) \\
\addlinespace[2pt]
Sacralised L5, four lumbar & L5 fused to the sacrum (Castellvi III or IV) and four lumbar labels; the fused body sits in the sacrum label: one box & 12 (9 / 3 / 0 / 0) & CTSP two-reader consensus; VerSe published read & 5 + 0 & 14 of 14, both sacralised rows & 8 (4) & 56.6 (5 CT); 51.0 (4 real) \\
\addlinespace[2pt]
Sacralised L5, five lumbar & L5 fused to the sacrum on a five-lumbar count; its own label (CTSP) or none (VerSe) & 16 (3 / 13 / 0 / 0) & as the row above & 6 + 0 & with the row above & none called & 61.0 (6 CT) \\
\addlinespace[2pt]
Castellvi I/II & enlarged or articulating transverse process; the count is unchanged & 31 (4 / 27 / 0 / 0) & CTSP two-reader consensus; VerSe published read & 12 + 0 & 15 of 15 & not screened & 75.8 (12 CT) \\
\addlinespace[2pt]
Lumbar rib & rib on L1; the count is unchanged & 14 (14 / 0 / 0 / 0) & CTSP release labels & 6 + 0 & 7 of 7 & not screened & 61.0 (6 CT) \\
\addlinespace[2pt]
Unassigned & four lumbar labels, source read normal; no subtype established & 4 (4 / 0 / 0 / 0) & CTSP release labels & 2 + 0 & 3 of 3 & -- & -- \\
\addlinespace[2pt]
L6 and T13 in one spine & none in any truth table & 0 & -- & -- & -- & -- & -- \\
\addlinespace
Any subtype & the rows above & 194 (52 / 135 / 4 / 3) & -- & 77 + 3 & 92 of 97 & 91 (7) & at-risk rows (L6, T13, four lumbar): 93.0 (51 CT); 64.6 (7 real) \\
\bottomrule
\end{tabular}
\par\vspace{1pt}
\begingroup\fontsize{7}{7.6}\selectfont\raggedright
CTSP = CTSpinoPelvic1K; TS = TotalSegmentator; C-Arm = C-ArmLSpine; L5, L6 = fifth, sixth lumbar vertebra; T13 = thirteenth thoracic vertebra.
\fnm{a}Distinct persons (patient key; two VerSe series pairs are one person each); rib convention: L6 = six rib-free vertebrae between the last rib-bearing vertebra and the sacrum, T13 = thirteen rib-bearing vertebrae. Counted 23 September 2026 from the subtype table; eight of twelve TotalSegmentator spines once called L6 from raw mask identifiers read normal on an independent rib count and are not counted.
\fnm{b}Whether a detector sees one body or two is decided by the mask: a fused body with its own label gives two boxes, one inside the sacrum label gives one, one left unlabelled gives none.
\fnm{c}Person-disjoint, subtype-stratified split, 50\,/\,10\,/\,40 train / validation / test within each subtype, every real transitional subject on the test side (the three C-ArmLSpine L6); built 23 September 2026, leak audit passed; Figure~\ref{fig:power}c draws these counts. The three real C-ArmLSpine L6 are scored only in cross-validation, as agreement with one labeller's read (Table~\ref{tab:naming}).
\fnm{d}Persons on the validation or test side of the built split who sit in the training split of the level bank behind the current encoder (the two most recent training sets); a model descending from those banks cannot score them, so a transitional test needs an arm trained without them.
\fnm{e}SPINEPS with VERIDAH on the planning CT deformably registered to the cone-beam CT, read by no person; on CTSP this labeller's L6 call matched the readers in 10 of 24 and its T13 call in 0 of 52. The screen also flagged three cases for re-reading. All calls sit on the test side; none counts as a transitional test person until verified. In brackets, counted 24 September 2026, the persons whom this labeller and an independent TotalSegmentator rib count, both requiring the thoracolumbar and lumbosacral junctions in the field, call transitional and who have an exported frame; no T13 can be counted, T1 lying outside nearly every field; four-lumbar persons are placed in the four-lumbar row.
\fnm{f}Wilson 95\,\% lower bound when every test person is named correctly, the best case at that count (``real'': the bracketed COBRA2026 persons, if each were verified); at risk = L6, T13 and the four-lumbar sacralised L5, the subtypes a count from the sacrum or from C2 misnames; an own-label sacralised L5, Castellvi I/II and a lumbar rib leave the count unchanged and are not in the at-risk bound; the pre-stated targets are 90\,\% for the exact level (35 persons with no error) and 85\,\% for variant sensitivity (22 carriers without a miss).
\par\endgroup
\end{minipage}}
\endgroup
\end{table}

% Table 3 (file name kept from the earlier layout; it prints third, after tab_corpus and tab4_transitional, since the
% LevelCheck table moved to paper 2 on 2026-09-24) --
% level naming by morphology, counting from the detected sacrum, and the coupled call, per stratum and
% field size. One measurement per row (uniform columns, no merged cells; a block heading sits in the first cell of
% its first row and repeated cells are left blank). Sources, one line per row, in tab3_naming.sources.md; the same
% rows in tab3_naming.md, which also keeps the rows cut from the page for space (five in field, AP three and one in
% field, the day-grouped AP and field rows, truth boxes at complete coverage, the within-one rows, the anchor-correct
% row, the DeepFluoro one-cluster count, the abstention rows, the coupled-minus-count, flag sensitivity, flag
% false-alarm and unflagged-exact rows, the coupling five-in-field rows and the oracle-box coupling block).
% Wilson cells (footnote f, formerly g): cells flagged ceiling_regime carry wilson_effective_n; REVIEW_STATS F3.
%   results/a4.json, a7.json         case-grouped five-fold CV and the abstention rule (k, n, clusters, percentile, BCa)
%   results/a9.json                  acquisition-day-grouped three-fold CV (numerators = round(exact x n))
%   results/a5.json                  plain film, patient-keyed validation, random window
%   results/counting.json            count from the detected sacrum (cells + pooled family)
%   results/coupling.json            classifier / count / coupled and the stop signal (percentile only)
% Transitional rows: the three C-ArmLSpine cases one labeller reads as L6 are shown as agreement with that read, not
% correctness (CLAIMS_LEDGER 6). Revision 5 (2026-09-24, REVISION_LOG3.md section 8): the row that owed the G6
% per-stratum result on held-out transitional persons moved to paper 2 (paper2_levelcheck/MOVED_FROM_PAPER1.md item 3);
% this table now prints as Table 3 (main.tex citation order 1 tab_corpus, 2 tab4_transitional, 3 tab3_naming).
% Added 2026-09-24 with Figure 4d (figures/figure4.record.json, panel d denominators): the L6 calls at confidence
% 0.74, from results/coupling.json#cells/det74~1kfull/calls_anchored (L6 4 of 233; non_ordinary_calls cases 3,
% every one flagged) and #l6_patients[arm=det74, field=full] (L6 on 1 of 4 anchored frames of the three patients).
% Revision 2 (2026-09-24, REVISION_LOG2.md; EVAL_LEAKS.md U1): the zero-shot rows are the lineage-complete v9 encoder
% and head (results/eval_gate/leaks/zeroshot_v9.json#zeroshot_v9/det/{LAT,AP}/kfull, eval_gate HELD OUT for all 148
% cases). Revision 2 round 2 (10:24 onward): every fold-model cell is now the v9-lineage re-run on the PRE truth from
% results/eval_gate/closure/cv_v9/cv_v9.json (families pre:cv5_det, pre:cvday_det: cells LAT|AP / kfull|k3|k1;
% Wilson where ceiling_regime is true) and closure/cv_v9/coupling_pre/coupling.json (rows det74/kfull, det/kfull,
% veridah_l6:det74/kfull, l6_patients_det74_full); every family eval_gate HELD OUT, overlap 0, unresolved 0. One
% measurement per row again (REVIEW2_RSNA Fix 8; the day-grouped field rows added). The zero-shot field row waits on
% the zs_v9_kfield job (grid, 10:43). The count-alone rows use no level encoder and are CLEAN (FIGREC:figure4/07-12).
% Page fit (build 40437080: 142 pt over the text width, notes clipped): the three plain-film selection-split rows and
% their note f moved to the supplement ("Plain-film naming on the selection split", every number and the caveat
% kept); the day-grouped one-in-field rows live in NUMBERS.md only (never quoted); notes re-lettered g->f, h->g. Footnote d adds the low-margin stop (REVIEW2_RSNA Fix 14; REVIEW2_STATS N4) and
% the classifier's lineage (N2); the plain-film caveat (closure 3c) now sits with its rows in the supplement.
% Revision 3 (2026-09-24, collect-and-submit-1; REVISION_LOG3.md): footnote d carries the abstention rule that left the
% Methods "The network" paragraph for room (posterior below 0.5 or decoded/single-crop disagreement; never tuned).
% Layout: landscape (rotated 23.2 cm minipage, see tab_corpus.tex); body 7 pt on 7.6 pt; the long descriptions of each
% test set live in the footnotes so that every row is one line.
\begin{table}[p]
\centering
\begingroup
\singlespacing
\fontsize{7}{7.6}\selectfont
\newcommand{\fnm}[1]{\textsuperscript{\fontsize{7}{7}\selectfont #1}}
\renewcommand{\arraystretch}{1.0}
\rotatebox[origin=c]{90}{\begin{minipage}{23.2cm}
\captionsetup{font=footnotesize,skip=4pt}
\centering
\caption{Level naming by morphology, counting from the detected sacrum, and the coupled call, per stratum and field size.}
\label{tab:naming}
% widths sum to 22.00 cm; 16 gaps of 1.9 pt (1.07 cm) -> 23.07 cm
\setlength{\tabcolsep}{1.9pt}
\begin{tabular}{%
>{\raggedright\arraybackslash}p{1.70cm}%
>{\raggedright\arraybackslash}p{3.40cm}%
>{\raggedright\arraybackslash}p{2.40cm}%
>{\raggedright\arraybackslash}p{5.00cm}%
>{\raggedright\arraybackslash}p{2.60cm}%
>{\raggedright\arraybackslash}p{2.10cm}%
>{\raggedright\arraybackslash}p{2.40cm}%
>{\raggedright\arraybackslash}p{2.40cm}}
\toprule
Mechanism & Test set; stratum; grouping & View; field\fnm{b}; operating point & Measure & \% ($k/n$) & Clusters (with an error) & 95\,\% percentile & BCa \\
\midrule
Morphology (level model) & C-ArmLSpine, agreed truth\fnm{a}; no C-ArmLSpine frame in training\fnm{g} & lateral; whole frame & exact level, full coverage & 96.4 (595/617) & 141 (6) & Wilson [92.3, 98.4]\fnm{f} & [92.5, 98.7] \\
 & & AP; whole frame & exact level & 43.7 (276/631) & 142 (83) & [34.9, 51.9] & [34.8, 51.8] \\
 & & lateral; three in field & exact level & 83.7 (970/1,159) & 142 (59) & [79.6, 87.5] & [79.4, 87.4] \\
 & & lateral; one in field & exact level & 63.8 (385/603) & 142 (118) & [59.8, 67.9] & [59.8, 67.7] \\
 & the same; case-grouped folds\fnm{g} & lateral; whole frame & exact level & 97.4 (601/617) & 141 (4) & Wilson [93.2, 99.0]\fnm{f} & [93.2, 99.4] \\
 & & AP; whole frame & exact level & 96.4 (608/631) & 142 (7) & Wilson [92.5, 98.3]\fnm{f} & [92.8, 98.4] \\
 & & lateral; three in field & exact level & 93.1 (1,079/1,159) & 142 (23) & [89.2, 96.2] & [88.5, 95.9] \\
 & & AP; three in field & exact level & 85.5 (1,005/1,176) & 143 (63) & [82.2, 88.6] & [82.1, 88.5] \\
 & & lateral; one in field & exact level & 74.8 (451/603) & 142 (94) & [70.9, 78.7] & [70.7, 78.5] \\
 & & AP; one in field & exact level & 67.5 (399/591) & 143 (113) & [64.2, 71.0] & [64.3, 71.1] \\
 & the same; day-grouped folds\fnm{g} & lateral; whole frame & exact level & 92.1 (568/617) & 141 (12) & [87.5, 96.0] & [86.7, 95.6] \\
 & & AP; whole frame & exact level & 94.8 (598/631) & 142 (9) & Wilson [90.1, 97.3]\fnm{f} & [90.2, 97.5] \\
 & & lateral; three in field & exact level & 87.7 (1,016/1,159) & 142 (38) & [83.5, 91.6] & [82.5, 91.1] \\
 & & AP; three in field & exact level & 86.0 (1,011/1,176) & 143 (63) & [82.6, 89.1] & [82.3, 88.9] \\
\addlinespace[1.5pt]
Count from the detected sacrum, no classifier\fnm{c} & C-ArmLSpine test, reported detector & whole frame; conf 0.74 & counting defined (a sacrum detected) & 82.1 (46/56) & 28 & [69.6, 92.9] & [69.6, 91.1] \\
 & & & ordinal correct, matched bodies & 95.4 (188/197) & 28 & [89.4, 99.5] & [85.2, 98.9] \\
 & & & truth bodies at the right ordinal, defined frames & 81.0 (188/232) & 28 & [73.8, 87.6] & [73.5, 87.1] \\
 & & whole frame; conf 0.25 & ordinal correct, matched bodies & 84.7 (199/235) & 28 & [77.4, 90.7] & [76.5, 90.4] \\
 & COBRA2026 test, held-out detector & conf 0.25 & ordinal correct, matched bodies & 71.1 (54/76) & 27 & [56.6, 80.3] & [57.7, 80.6] \\
 & pooled over the three held-out devices & conf 0.25 & ordinal correct, matched bodies & 81.0 (306/378) & 56 & [75.7, 85.8] & [75.7, 85.8] \\
\addlinespace[1.5pt]
Coupled: classifier and count\fnm{d} & C-ArmLSpine anchored frames, agreed truth & whole frame; conf 0.74 & count alone, counted bodies & 97.0 (979/1,009) & 130 & [94.2, 99.2] & -- \\
 & & & classifier alone, counted bodies & 97.7 (986/1,009) & 130 & [95.9, 99.2] & -- \\
 & & & coupled, counted bodies & 98.2 (991/1,009) & 130 & [96.5, 99.5] & -- \\
 & & & coupled minus count; minus classifier (points) & $+$1.2; $+$0.5 & 130 & [0.0, 3.2]; [$-$1.5, 2.4] & -- \\
 & & & disagreement; stop (13 disagreement, 10 low margin); erring caught & 5.6 (13/233); 9.9 (23/233); 5/6 & 132 & [2.6, 9.2]; [5.7, 14.5] & -- \\
 & & & L6 calls on five-lumbar frames, every one stopped & 3/233 & 2 & no interval & -- \\
 & & whole frame; conf 0.25 & count alone, counted bodies & 96.1 (1,104/1,149) & 139 & [93.5, 98.3] & -- \\
 & & & classifier alone; coupled, counted bodies & 97.7 (1,122/1,149); 97.1 (1,116/1,149) & 139 & [95.8, 99.2]; [95.3, 98.8] & -- \\
 & & & disagreement frames & 7.1 (19/266) & 142 & [3.8, 10.8] & -- \\
\addlinespace[1.5pt]
Transitional, real exposures & three C-ArmLSpine cases read as L6 by one labeller\fnm{e} (Table~\ref{tab:transitional}) & whole frame; conf 0.25 and 0.74 & agreement with that read on the L6 body (coupled; classifier); anchored frames called L6 (conf 0.74) & 1/5; 0/5; 1 of 4 & 3 & no interval & -- \\
\bottomrule
\end{tabular}
\par\vspace{1pt}
\begingroup\fontsize{7}{7.6}\selectfont\raggedright
AP = anteroposterior; BCa = bias-corrected and accelerated bootstrap interval; conf = detector confidence threshold. Intervals are 95\,\% bootstraps of 2,000 resamples clustered by case (intraoperative) or patient (film); ``--'' = not computed.
\fnm{a}Agreed truth = the levels both automatic CBCT labellers name (814 of 910 vertebrae; T11 and most of T12 withheld); ordinary spines. Case-grouped = five folds by case, day-grouped = three folds by acquisition day (a case is an acquisition session, not a proven person); ``with an error'' = clusters carrying at least one error.
\fnm{b}Field: whole frame, or crops holding three or one vertebra. Exact level scores every vertebra the film-trained keypoint detector boxed (confidence 0.35; 617 of 676 agreed lateral; with the 59 undetected counted wrong, 88.0\,\% zero-shot) on its decoded name at full coverage; naming matches a detection to the agreed vertebral-body box containing its centre (overlap at least 0.3, one to one).
\fnm{c}Ordinal = rank of the detected body away from the detected sacrum, scored name-free against the rank from the truth sacrum; defined when truth has a sacrum and one was detected; one-to-one matching at intersection over union $\ge$ 0.5, whole-vertebra convention (supplement, Table~S2); a missed truth body counts as wrong in the defined-frames row. C-ArmLSpine test = 28 cases, 56 AP and lateral frames, device never seen; COBRA2026 test = 27 patients, 609 frames with a sacrum, detector trained on DeepFluoro and C-ArmLSpine real frames; pooled = each real-frame detector on its unseen device.
\fnm{d}C-ArmLSpine five-fold CV frames with agreed truth (145 cases, 290 frames), the sacrum from the held-out detector. Classifier = the level model's name (the case-grouped fine-tuned fold models; AP and lateral frames pooled); count = the ordinary alphabet at the ordinal (L5 at one), no classifier; coupled = the alphabet (ordinary, L6, T13, both) that best fits the classifier's posteriors at the ordinals; the read stops when classifier and call disagree or the winning alphabet's margin over the runner-up is below the fifth percentile of correct frames in the other four folds; the level model abstains on a body when its decoded and single-crop names disagree or its posterior is below 0.5, a default never tuned; percentile intervals only.
\fnm{e}Agreement with that read, not correctness: the two labellers segment the same five bodies and disagree only on the names (L1--L5 against L2--L6) on a lumbar-only CBCT; six frames, three clusters, no interval.
\fnm{f}Wilson 95\,\% interval on the design-effect-adjusted count where the result file flags errors in few clusters.
\fnm{g}Zero-shot = the encoder and its recurrent head with no C-ArmLSpine case in training (checked for all 148 cases); the fold models are fine-tuned from that encoder, every fold held out by list intersection.
\par\endgroup
\end{minipage}}
\endgroup
\end{table}


\end{document}
